
%\needspace{4\baselineskip}
\color{uniblue}{\subsection{Автоматическое ограничение снижения напряжения (АОСН)}}
\color{black}

\begin{enumerate}[label=\arabic{section}.\arabic{subsection}.\arabic{enumi}, labelsep=4pt, leftmargin=0pt, itemindent=57pt, itemsep=0pt, parsep=5pt] % если что можно крутить itemsep
%%%%%%%%%%%%%%%%%%%%%%%%%%%%%%%%%%%%%%%%%%%%%%%%%%%%%%%%%%%  ОБЩИЕ СВЕДЕНИЯ %%%%%%%%%%%%%%%%%%%%%%%%%%%%%%%%%%%%%%%%%%%%%%%%

\item Общие сведения

\begin{enumerate}[label=\arabic{section}.\arabic{subsection}.\arabic{enumi}.\arabic*, labelsep=4pt, leftmargin=0em, itemindent=65pt, parsep=0pt]

\item
Автоматическое ограничение снижения напряжения необходимо для предотвращения развития аварий, сопровождающихся глубоким снижением напряжения. Функция применяется для предотвращения снижения напряжения в узлах энергосистемы, вызванного дефицитом реактивной мощности до значения, опасного по условиям устойчивости энергосистемы и возникновения лавины напряжения, путем отключения части нагрузки с последующим включением отключенной нагрузки при восстановлении напряжения.

\item
АОСН включает в себя следующие функции:
\begin{itemize}
    \item три очереди АОСН (АОСН1, АОСН2, АОСН3);
    \item 15 функций для реализации управляющих воздействий очередей АОСН (УВ1 АОСН ... УВ15 АОСН);
    \item три функции АПВН (АПВН1, АПВН2, АПВН3);
    \item 15 функций для реализации управляющих воздействий функций АПВН (УВ1 АПВН ... УВ15 АПВН);
    \item общая функция объединения срабатывания функций в составе функционального блока <<АОСН>>.   
\end{itemize}

\item
Основные характеристики АОСН:
\begin{itemize}
    \item основная погрешность срабатывания ИО частоты не более 20 мГц;
    \item основная погрешность возврата ИО частоты не более 20 мГц;
    \item основная погрешность ИО частоты в рабочем диапазоне при напряжении входного сигнала 100 В не более 10 мГц;
    \item время срабатывания и возврата ИО частоты не более 60 мс;
    \item основная относительная приведенная погрешность измерения по напряжению в рабочем диапазоне при частоте входного сигнала (50 ± 0,5) Гц составляет ± 3,0\%;
    \item дополнительная погрешность измерения напряжения и тока при изменении частоты входного сигнала в рабочем диапазоне на каждый 1 Гц составляет ± 0,5\%.
\end{itemize}

\item
Функциональный блок <<АОСН>> может быть полностью выведен из работы оперативно путем активации сигнала <<ОВ АОСН/АПВН>>.

\end{enumerate}

%%%%%%%%%%%%%%%%%%%%%%%%%%%%%%%%%%%%%%%%%%%%%%%%%%%%%%%%%  ОЧЕРЕДЬ АОСН %%%%%%%%%%%%%%%%%%%%%%%%%%%%%%%%%%%%%%%%%%%%%%%%

\item Очереди АОСН (АОСН1, АОСН2, АОСН3)

\begin{enumerate}[label=\arabic{section}.\arabic{subsection}.\arabic{enumi}.\arabic*, labelsep=4pt, leftmargin=0em, itemindent=65pt, parsep=0pt]

\item
Ввод каждой очереди АОСН в работу на этапе параметрирования устройства осуществляется своим программным переключателем SGF1 <<Ввод\_функции>>. Информация о введенной очереди отображается выходным сигналом для каждой очереди <<АОСН(n): Ввод>>, здесь и далее в этом подразделе: (n) -- номер очереди, $n\in\{1,2,3\}$.

\item
Очереди <<АОСН>> могут быть выведены из работы оперативно путем активации общего сигнала вывода всех трех очередей <<ОВ АОСН>> либо по отдельности каждая своим сигналом оперативного вывода(<<ОВ АОСН(n)>>). Состояние очереди характеризуется активным состоянием сигнала <<АОСН(n): Оперативный вывод>> на выходе функции.

\item
При активном состоянии входного сигнала <<АОСН(n) на сигн>>  действие очереди (n) АОСН переводится на сигнал.

\item
В каждой из очередей предусмотрены следующие измерительные органы:
\begin{itemize}
    \item ИО минимального действия по значению линейным напряжениям. Срабатывание осуществляется при снижении минимального из поданных напряжений ниже уставки;
    \item два ИО минимального действия по напряжению ПП на двух секциях. Срабатывание происходит при снижении напряжений прямой последовательности на каждой из двух секциях ниже уставок;
    \item ИО максимального действия по напряжению обратной последовательности. Срабатывание происходит при превышении значения напряжения уставки напряжения ОП;
    \item ИО по снижению частоты. Функция реагирует на снижение значения частоты ниже уставки;
    \item ИО по скорости снижения напряжения $dU/dt$. Срабатывание происходит при превышении скорости снижения напряжения уставки.
\end{itemize}

\item
Пусковым условием АОСН является одновременное выполнение условий:
\begin{itemize}
    \item отсутствие сигналов: срабатывания напряжения обратной последовательности, скорости снижения напряжения, линейных напряжений (реализованных через выдержку времени на срабатывание T3 и подхвата T4), снижение частоты ниже порогового значения;
    \item снижение напряжения на обеих секциях;
    \item отсутствие срабатываний БНН и блокировки АОСН.
\end{itemize}

\item
В случае снижения напряжения прямой последовательности основной и смежной секции до порога срабатывания очереди АОСН, запускается выдержка времени данной очереди и в случае сохранения пусковых условий, происходит срабатывание данной очереди АОСН. Кроме того, необходимо отсутствие остальных блокирующих сигналов АОСН, как общих, так и по каждой очереди.

\item
Блокировка пуска очереди АОСН осуществляется при выполнении одного из условий:
\begin{itemize}
    \item срабатывание ИО ОП при введенной уставке SGF2;
    \item срабатывание ИО по скорости снижения напряжения при введенной уставке SGF3;
    \item срабатывание ИО по скорости снижения напряжения по истечению времени таймера Т3. При пропадании сигнала, блокировка снимается по истечению времени таймера Т4;
    \item наличие входных сигналов <<БНН: Срабатывание>>, <<Блокировка АОСН(n)>>.
\end{itemize}

\item
Срабатывание очереди АОСН происходит при наличии пускового сигнала через выдержку времени T1 и при отсутствии сигнала <<АОСН(n) на сигн>>. При наличии сигнала <<АОСН(n) на сигн>> действие очереди переводится на сигнал <<АОСН(n): Срабатывание сигн>>, при этом выходной сигнал <<АОСН(n): Срабатывание>> блокируется.
\item
Для исключения кратковременных пусков ступеней, в соответствии с ГОСТ Р 70411-2022, предусмотрена возможность продления сигнала пуска ИО частоты на время, задаваемое уставкой Т2.

\item
Алгоритм работы очереди АОСН выполнен в соответствии с рисунком \ref{aosn:img1}. В таблице \ref{aosn:tbl1} приведены параметры, необходимые для настройки очереди АОСН.

\small
\begin{longtable}{|>{\centering\arraybackslash}m{5.3cm}|>{\centering\arraybackslash}m{3.3cm}|>{\centering\arraybackslash}m{4.2cm}|>{\centering\arraybackslash}m{1.8cm}|>{\centering\arraybackslash}m{1cm}|}
\caption{Параметры для настройки очереди АОСН\hfill\vspace{-0.5\baselineskip}}\label{aosn:tbl1}\\ 
\hline
\rowcolor{gray!30}
Параметр (Параметр на ИЧМ) & Условное обозначение на схеме & Значение/ Диапазон & Единица измерения & Шаг \\ 
\hline
\endfirsthead
\caption*{\hspace{3pt}\emph{Продолжение таблицы \ref{aosn:tbl1}\hfill\vspace{-0.5\baselineskip}}} \\ % сделано по ГОСТ 2.105 п.6.8.7
\hline
\rowcolor{gray!30}
Параметр (Параметр на ИЧМ) & Условное обозначение на схеме & Значение/ Диапазон & Единица измерения & Шаг \\ 
%\hline
\endhead
%\hline
\endfoot
%\hline
\endlastfoot
%==+t1*ZSPTOC1|LVZSTOC
\centering Ввод функции в работу (Ввод\_функции) & \centering SGF1 & \centering 0 = Не предусмотрено\\1 = Предусмотрено & \centering -- & \centering \arraybackslash -- \\
\hline
\centering Режим оперативного блокирования ЧС (Реж\_БЧС) & \centering SGF8 & \centering 0 = Не предусмотрено\\1 = Предусмотрено & \centering -- & \centering \arraybackslash -- \\
\hline
\centering Режим контроля от БНТ (Реж\_БНТ) & \centering SGF7 & \centering 0 = Не предусмотрено\\1 = Предусмотрено & \centering -- & \centering \arraybackslash -- \\
\hline
\centering Ток срабатывания НП (3I0ср) & \centering 3I0$>$ & \centering 0,10 ... 30,00 & \centering о.е. & \centering \arraybackslash 0,01 \\
\hline
\centering Вывод направленности при АУ (Ненапр\_АУ) & \centering SGF4 & \centering 0 = Не предусмотрено\\1 = Предусмотрено & \centering -- & \centering \arraybackslash -- \\
\hline
\centering Режим направленности (Направленность) & \centering SGF2 & \centering 1 = Ненаправленная\\2 = Прямонаправленная\\3 = Обратнонаправленная & \centering -- & \centering \arraybackslash -- \\
\hline
\centering Режим ОНМ НП при неисправности ЦН (Реж\_БНН\_ОНМ\_НП) & \centering SGF3 & \centering 0 = Блокировка\\1 = Деблокировка & \centering -- & \centering \arraybackslash -- \\
\hline
\centering Режим ПО 3U0 при неисправности ЦН (Реж\_БНН\_3U0) & \centering SGF5 & \centering 0 = Блокировка\\1 = Деблокировка & \centering -- & \centering \arraybackslash -- \\
\hline
\centering Напряжение срабатывания НП (3U0ср) & \centering 3U0$>$ & \centering 2,0 ... 100,0 & \centering В & \centering \arraybackslash 0,1 \\
\hline
\centering Режим контроля от ПО 3U0 (Реж\_3U0) & \centering SGF6 & \centering 0 = Не предусмотрено\\1 = Предусмотрено & \centering -- & \centering \arraybackslash -- \\
\hline
\centering Характеристика срабатывания (Характеристика) & \centering -- & \centering 16 = Независимая\\11 = Кривая A\\10 = Кривая B\\12 = Кривая C\\4 = Кривая D\\2 = Кривая E\\1 = Кривая F\\17 = Кривая RXIDG & \centering -- & \centering \arraybackslash -- \\
\hline
\centering Коэффициент регулировки времени срабатывания зависимой характеристики (К) & \centering -- & \centering 0,05 ... 15,00 & \centering -- & \centering \arraybackslash 0,01 \\
\hline
\centering Независимая выдержка времени срабатывания (Tср) & \centering T1 & \centering 0,00 ... 100,00 & \centering с & \centering \arraybackslash 0,01 \\
\hline
%===t1
\end{longtable}
\normalsize

\vspace{3mm}
\begin{figure}[H]
\centering
\includegraphics[width=0.95\textwidth,height=0.95\textheight,keepaspectratio]{\fbpath/205. АОСН/_latex/img/aosn.pdf}
\captionof{figure}{Функциональная схема очереди АОСН}\label{aosn:img1}
\end{figure}

\end{enumerate}

%%%%%%%%%%%%%%%%%%%%%%%%%%%%%%%%%%%%%%%%%%%%%%%%%%%%%%%%% УВ АОСН %%%%%%%%%%%%%%%%%%%%%%%%%%%%%%%%%%%%%%%%%%%%%%%%

\item Реализация управляющих воздействий от очередей АОСН (УВ АОСН)

\begin{enumerate}[label=\arabic{section}.\arabic{subsection}.\arabic{enumi}.\arabic*, labelsep=4pt, leftmargin=0em, itemindent=65pt, parsep=0pt]

\item
Для реализации управляющих воздействий (например, на отключение нагрузки) при срабатывании очередей АОСН предусмотрено 15 индивидуально параметрируемых функций: УВ1 АОСН ... УВ15 АОСН (далее в этом подразделе <<УВ(m)~АОСН>>, где (m) -- номер функции УВ, $\{m \in \mathbb{N} \mid 1 \leq m \leq 15\}$.  

\item
Уставкой SGF2 контролируется от какой ступени АОСН будут формироваться УВ. Срабатывание функции происходит по истечении выдержки времени таймера Т1, наличии сигнала пуска соответствующей ступени и отсутствии сигнала <<УВ(m) АОСН на сигн>>. Наличие сигнала «УВ(m)АОСН на сигн» переводит действие функции на сигнал, при этом сигнал срабатывания функции будет заблокирован.






\item
Функция <<ОНМ НП>> автоматически вводится в работу, если введена в работу одна из ступеней ТЗНП и в этой ступени настроен контроль направленности.

\item
В качестве воздействующих величин ОНМ НП использует значения первой гармоники тока нулевой последовательности и напряжения нулевой последовательности.

\item
Направленный элемент срабатывает при выполнении следующих условий:
\begin{itemize}
    \item напряжение НП превышает 1 вольт;
    \item ток НП превышает значение $0,04 \cdot I_{\textsl{ном}}$ НП;
    \item вектор полной мощности нулевой последовательности, определяемый по углу током НП и напряжением НП, находится в области срабатывания (прямое или обратное направление).
\end{itemize}

\item
Для определения области срабатывания задается уставка $\varphi_{\textsl{м.ч.}}$. Отсчет угла максимальной чувствительности производится от вектора напряжения НП и считается положительным, если направление отсчета совпадает с направлением вращения против часовой стрелки. Угол сектора отсчитывается от угла максимальной чувствительности в обе стороны и составляет 85°.

\item 
Сигнал <<ОНМ НП: Прямое>> появляется, если вектор полной мощности НП находится в диапазоне от $\varphi_{\textsl{м.ч.}}$ - 85° до  $\varphi_{\textsl{м.ч.}}$ + 85°. Сигнал <<ОНМ НП: Обратное>> появляется, если вектор полной мощности НП находится в диапазоне от ($\varphi_{\textsl{м.ч.}}$ - 180°) + 85° до ($\varphi_{\textsl{м.ч.}}$ - 180°) - 85°. Если вектор полной мощности НП находится в зоне <<нечувствительности>>, то ни один из вышеуказанных сигналов не появляется.

\item 
Поясняющая диаграмма определения направления мощности приведена на рисунке \ref{tznp:imgonm0}.

\vspace{3mm}
\begin{figure}[H]
\centering
\includegraphics[width=0.8\textwidth,height=0.8\textheight,keepaspectratio]{\fbpath/106. ТЗНП ОЛ/_latex/img/char.pdf}
\captionof{figure}{Характеристика органа направления мощности нулевой последователности}\label{tznp:imgonm0}
\end{figure}

\item
Алгоритм работы функции <<ОНМ НП>> выполнен в соответствии с рисунком \ref{tznp:imgonm1}. В таблице \ref{tznp:tblonm} приведены параметры, необходимые для настройки функции <<ОНМ НП>>.

\small
\begin{longtable}{|>{\centering\arraybackslash}m{5.3cm}|>{\centering\arraybackslash}m{3.3cm}|>{\centering\arraybackslash}m{4.2cm}|>{\centering\arraybackslash}m{1.8cm}|>{\centering\arraybackslash}m{1cm}|}
\caption{Параметры для настройки функции <<ОНМ НП>>\hfill\vspace{-0.5\baselineskip}}\label{tznp:tblonm}\\ 
\hline
\rowcolor{gray!30}
Параметр (Параметр на ИЧМ) & Условное обозначение на схеме & Значение/ Диапазон & Единица измерения & Шаг \\ 
\hline
\endfirsthead
\caption*{\hspace{3pt}\emph{Продолжение таблицы \ref{tznp:tblonm}\hfill\vspace{-0.5\baselineskip}}} \\ % сделано по ГОСТ 2.105 п.6.8.7
\hline
\rowcolor{gray!30}
Параметр (Параметр на ИЧМ) & Условное обозначение на схеме & Значение/ Диапазон & Единица измерения & Шаг \\ 
%\hline
\endhead
%\hline
\endfoot
%\hline
\endlastfoot
%==+t1*RDIR1|LVZSTOC
\centering Угол максимальной чувствительности (Фмч) & \centering φмч & \centering -179,0 ... 180,0 & \centering град. & \centering \arraybackslash 0,1 \\
\hline
%===t1
\end{longtable}
\normalsize

\vspace{3mm}
\begin{figure}[H]
\centering
\includegraphics[width=0.9\textwidth,height=0.9\textheight,keepaspectratio]{\fbpath/106. ТЗНП ОЛ/_latex/img/onmnp.pdf}
\captionof{figure}{Функциональная схема <<ОНМ НП>>}\label{tznp:imgonm1}
\end{figure}

\end{enumerate}

%%%%%%%%%%%%%%%%%%%%%%%%%%%%%%%%%%%%%%%%%%%%%%%%%%%%%%%%% ОРГАН БНТ %%%%%%%%%%%%%%%%%%%%%%%%%%%%%%%%%%%%%%%%%%%%%%%%

\item Блокировка при броске намагничивающего тока НП (БНТ НП)

\begin{enumerate}[label=\arabic{section}.\arabic{subsection}.\arabic{enumi}.\arabic*, labelsep=4pt, leftmargin=0em, itemindent=65pt, parsep=0pt]

\item
Функция блокировки при броске намагничивающего тока НП (<<БНТ НП>>) позволяет заблокировать защиту нулевой последовательности при броске тока намагничивания НП, возникающем при постановке силового трансформатора под напряжение. 
\item
Ввод функции <<БНТ НП>> осуществляется автоматически при активации режим контроля броска тока намагничивания введенной в работу ступени ТЗНП.
\item
Условие пуска выполняется при одновременном наличии следующих условий:
\begin{itemize}
\item отношение действующего значения тока НП второй гармоники (\textsl{3I0 2г}) к действующему значению тока НП первой гармоники (\textsl{3I0}) превышает уставку <<Ih2/Ih1>>;
\item отсутствует срабатывание порогового элемента, реагирующего на превышение действующего значения первой гармоники тока НП (\textsl{3I0}) над уставкой <<Iср>>;
\item сформирован пуск одной из ступеней ТЗНП.
\end{itemize}

\item
Алгоритм работы функции <<БНТ НП>> выполнен в соответствии с рисунком \ref{tznp:img3}. В таблице \ref{tznp:tbl3} приведены параметры, необходимые для настройки функции <<БНТ НП>>.

\vspace{3mm}
\begin{figure}[H]
\centering
\includegraphics[width=0.8\textwidth,height=0.8\textheight,keepaspectratio]{\fbpath/106. ТЗНП ОЛ/_latex/img/bntnp.pdf}
\captionof{figure}{Функциональная схема <<БНТ НП>>}\label{tznp:img3}
\end{figure}

\small
\begin{longtable}{|>{\centering\arraybackslash}m{5.3cm}|>{\centering\arraybackslash}m{3.3cm}|>{\centering\arraybackslash}m{4.2cm}|>{\centering\arraybackslash}m{1.8cm}|>{\centering\arraybackslash}m{1cm}|}
\caption{Параметры для настройки функции <<БНТ НП>>\hfill\vspace{-0.5\baselineskip}}\label{tznp:tbl3}\\ 
\hline
\rowcolor{gray!30}
Параметр (Параметр на ИЧМ) & Условное обозначение на схеме & Значение/ Диапазон & Единица измерения & Шаг \\ 
\hline
\endfirsthead
\caption*{\hspace{3pt}\emph{Продолжение таблицы \ref{tznp:tbl3}\hfill\vspace{-0.5\baselineskip}}} \\ % сделано по ГОСТ 2.105 п.6.8.7
\hline
\rowcolor{gray!30}
Параметр (Параметр на ИЧМ) & Условное обозначение на схеме & Значение/ Диапазон & Единица измерения & Шаг \\ 
%\hline
\endhead
%\hline
\endfoot
%\hline
\endlastfoot
%==+t1*PHAR1|LVZSTOC
\centering Отношение тока второй гармоники к току основной гармоники (Ih2/Ih1) & \centering I2г/I1г$>$ & \centering 5 ... 100 & \centering \% & \centering \arraybackslash 1 \\
\hline
\centering Ток блокировки НП (3I0ср) & \centering 3I0$>$ & \centering 0,1 ... 15,0 & \centering о.е. & \centering \arraybackslash 0,1 \\
\hline
%===t1
\end{longtable}
\normalsize

\end{enumerate}

%%%%%%%%%%%%%%%%%%%%%%%%%%%%%%%%%%%%%%%%%%%%%%%%%%%%%%%%% ОРГАН АУ %%%%%%%%%%%%%%%%%%%%%%%%%%%%%%%%%%%%%%%%%%%%%%%%

\item Автоматическое ускорение ТЗНП (АУ ТЗНП)

\begin{enumerate}[label=\arabic{section}.\arabic{subsection}.\arabic{enumi}.\arabic*, labelsep=4pt, leftmargin=0em, itemindent=65pt, parsep=0pt]

\item
В устройстве предусмотрен режим автоматического ускорения ступеней ТЗНП. Выбор ускоряемых ступеней защиты осуществляется из соображений надежного охвата, защищаемого энергообъекта.
\item
Предусмотрена возможность автоматического ускорения для любой ступени ТЗНП. Ускоряемая ступень выбирается уставкой SGF1. Автоматическое ускорение срабатывает без выдержки времени.

\item
Алгоритм работы <<АУ ТЗНП>> выполнен в соответствии с рисунком \ref{tznp:imgau}. В таблице \ref{tznp:tblau} приведены параметры, необходимые для настройки функции <<АУ ТЗНП>>.

\begin{figure}[H]
\centering


\includegraphics[width=0.6\textwidth,height=0.6\textheight,keepaspectratio]{\fbpath/106. ТЗНП ОЛ/_latex/img/au.pdf}
\captionof{figure}{Функциональная схема АУ ТЗНП}\label{tznp:imgau}
\vspace{3mm}
\end{figure}

\small
\begin{longtable}{|>{\centering\arraybackslash}m{5.3cm}|>{\centering\arraybackslash}m{3.3cm}|>{\centering\arraybackslash}m{4.2cm}|>{\centering\arraybackslash}m{1.8cm}|>{\centering\arraybackslash}m{1cm}|}
\caption{Параметры для настройки функции <<АУ ТЗНП>>\hfill\vspace{-0.5\baselineskip}}\label{tznp:tblau}\\ 
\hline
\rowcolor{gray!30}
Параметр (Параметр на ИЧМ) & Условное обозначение на схеме & Значение/ Диапазон & Единица измерения & Шаг \\ 
\hline
\endfirsthead
\caption*{\hspace{3pt}\emph{Продолжение таблицы \ref{tznp:tblau}\hfill\vspace{-0.5\baselineskip}}} \\ % сделано по ГОСТ 2.105 п.6.8.7
\hline
\rowcolor{gray!30}
Параметр (Параметр на ИЧМ) & Условное обозначение на схеме & Значение/ Диапазон & Единица измерения & Шаг \\ 
%\hline
\endhead
%\hline
\endfoot
%\hline
\endlastfoot
%==+t1*AUAPSSU1|LVZSTOC
\centering Выбор ускоряемой ступени (Выбор\_ст) & \centering SGF1 & \centering 1 = 1 ступень\\2 = 2 ступень & \centering -- & \centering \arraybackslash -- \\
\hline
%===t1
\end{longtable}
\normalsize

\end{enumerate}
%%%%%%%%%%%%%%%%%%%%%%%%%%%%%%%%%%%%%%%%%%%%%%%%%%%%%%%%% КОНЕЦ ОРГАН АУ %%%%%%%%%%%%%%%%%%%%%%%%%%%%%%%%%%%%%%%%%%%%%%%%

%%%%%%%%%%%%%%%%%%%%%%%%%%%%%%%%%%%%%%%%%%%%%%%%%%%%%%%%% ОРГАН ОУ %%%%%%%%%%%%%%%%%%%%%%%%%%%%%%%%%%%%%%%%%%%%%%%%

\item Оперативное ускорение ТЗНП (ОУ ТЗНП)

\begin{enumerate}[label=\arabic{section}.\arabic{subsection}.\arabic{enumi}.\arabic*, labelsep=4pt, leftmargin=0em, itemindent=65pt, parsep=0pt]

\item
Оперативное ускорение предусматривает ускорение любой ступени ТЗНП. Оперативное ускорение может вводиться:
\begin{itemize}
    \item при выводе основной защиты оборудования;
    \item при опробовании оборудования рабочим напряжением;
    \item при операциях с разъединителями на противоположном конце ВЛ.
\end{itemize}

\item
В данном алгоритме при введенном ключе <<Ввод ОУ ТЗНП>> предусмотрена возможность оперативного ускорения любой ступени ТЗНП. Выбор ускоряемой ступени осуществляется программным переключателем SGF1 <<Выбор\_ст\_ОУ>>. При введенной функции и появлению сигнала пуска выбранной ступени по истечению выдержки времени, определяемой уставкой T1, формируется сигнал срабатывания.

\item
Алгоритм работы <<ОУ ТЗНП>> выполнен в соответствии с рисунком \ref{tznp:imgou}. В таблице \ref{tznp:tblou} приведены параметры, необходимые для настройки функции <<ОУ ТЗНП>>.

\begin{figure}[H]
\centering

\includegraphics[width=0.8\textwidth,height=0.8\textheight,keepaspectratio]{\fbpath/106. ТЗНП ОЛ/_latex/img/ou.pdf}
\captionof{figure}{Функциональная схема ОУ ТЗНП}\label{tznp:imgou}
\vspace{3mm}
\end{figure}

\small
\begin{longtable}{|>{\centering\arraybackslash}m{5.3cm}|>{\centering\arraybackslash}m{3.3cm}|>{\centering\arraybackslash}m{4.2cm}|>{\centering\arraybackslash}m{1.8cm}|>{\centering\arraybackslash}m{1cm}|}
\caption{Параметры для настройки функции <<ОУ ТЗНП>>\hfill\vspace{-0.5\baselineskip}}\label{tznp:tblou}\\ 
\hline
\rowcolor{gray!30}
Параметр (Параметр на ИЧМ) & Условное обозначение на схеме & Значение/ Диапазон & Единица измерения & Шаг \\ 
\hline
\endfirsthead
\caption*{\hspace{3pt}\emph{Продолжение таблицы \ref{tznp:tblou}\hfill\vspace{-0.5\baselineskip}}} \\ % сделано по ГОСТ 2.105 п.6.8.7
\hline
\rowcolor{gray!30}
Параметр (Параметр на ИЧМ) & Условное обозначение на схеме & Значение/ Диапазон & Единица измерения & Шаг \\ 
%\hline
\endhead
%\hline
\endfoot
%\hline
\endlastfoot
%==+t1*OPAPSSU1|LVZSTOC
\centering Выбор оперативно ускоряемой ступени (Выбор\_ст\_ОУ) & \centering SGF1 & \centering 1 = 1 ступень\\2 = 2 ступень & \centering -- & \centering \arraybackslash -- \\
\hline
\centering Выдержка времени срабатывания оперативного ускорения (T\_ОУ) & \centering T1 & \centering 0,00 ... 3,00 & \centering с & \centering \arraybackslash 0,01 \\
\hline
%===t1
\end{longtable}
\normalsize

\end{enumerate}
%%%%%%%%%%%%%%%%%%%%%%%%%%%%%%%%%%%%%%%%%%%%%%%%%%%%%%%%% КОНЕЦ ОРГАН ОУ %%%%%%%%%%%%%%%%%%%%%%%%%%%%%%%%%%%%%%%%%%%%%%%%

\end{enumerate}
